\uuid{cfa8}
\titre{Matrice/Somme/Produit}
\auteur{hanani ; mzari}
\organisation{exo7}
\licence{CC BY-NC-SA 4.0}
\datecreate{2020-01-20}
\niveau{L1}
\module{Algèbre}
\chapitre{Matrice}
\sousChapitre{Propriétés élémentaires, généralités}
\difficulte{2}
\theme{Calcul matriciel, Matrice, Somme, Produit}
\exo7qcmid{lille-2.matrices.007}

\contenu{
\begin{question}
\qclassification{108.01}
\qtags{motcle=Matrice/Somme/Produit}
Soit $A,B$ et $C$ des matrices d'ordre $n\ge 1$. Quelles sont les assertions vraies ?
\begin{answers}
\bad{$AB=0 \Rightarrow A=0 \, \mbox{ou} \,  B=0$}.
\bad{$A(BC) =(AC)B$}.
\good{$A(B+C)=AC+AB$}.
\bad{$(A+B)^2=A^2+2AB+B^2$}.
\end{answers}
\begin{explanations} Le produit de deux matrices peut \^etre nul sans que l'une des deux matrices soit nul. Contre-exemple : avec $A=
\left(\begin{array}{rc}0&0\\
1&0\end{array}\right)$ on a : $A^2=0$. Le produit de deux matrices est associatif et distributif par rapport à l'addition. Comme le produit n'est pas commutatif, en général, $(A+B)^2=A^2+AB+BA+B^2 \neq A^2+2AB+B^2$.
\end{explanations}
\end{question}
}
